\documentclass[10pt,a4paper]{amsart}
\usepackage[margin=19mm]{geometry}
\usepackage{amsmath,amssymb,mathtools}
\usepackage[scr=rsfs,cal=euler]{mathalfa}
\usepackage{xfrac}
\usepackage[unicode,hidelinks,bookmarks=false]{hyperref}
\newcommand{\ie}{i.e.\ }
\newcommand{\cf}{cf.\ }
\newcommand{\resp}{respectively\ }
\newcommand{\Subsection}{\subsection}
\providecommand{\nobreakdash}{\nobreak\hbox{-}}
\setlength{\emergencystretch}{2em}
\multlinegap0pt
\usepackage{bm}
\usepackage{bbm}
\usepackage{dsfont}
\usepackage{xspace}
\usepackage{cancel}
\usepackage{mathtools}
\usepackage{amsthm}
\makeatletter
\@ifundefined{theorem}{\newtheorem{theorem}{Theorem}}{}
\@ifundefined{lem}{\newtheorem{lem}[theorem]{Lemma}}{}
\makeatother
\def\Cl{\mathop\mathrm{Cl}}
\def\N{{\mathbb N}}
\def\R{{\mathbb R}}
\def\eps{\varepsilon}
\newcommand{\C}{\mathbb{C}}
\newcommand{\im}{\operatorname{Im}}
\newcommand{\re}{\operatorname{Re}}
\renewcommand{\phi}{\varphi}
\title[The zero entropy locus for the Lozi maps]{The zero entropy locus for the Lozi maps}
\author{M. Misiurewicz, S. \v{S}timac}
\date{Source: arXiv:2411.17836v2 (CC BY 4.0)}
\begin{document}
\maketitle

\noindent\emph{Source and licence.} The zero entropy locus for the Lozi maps. Version arXiv:2411.17836v2 (CC BY 4.0), distributed under the Creative Commons Attribution 4.0 International licence, as recorded in the arXiv licence metadata for this exact version and shown on the abstract page https://arxiv.org/abs/2411.17836v2. Full rights notice: licenses/arxiv-2411.17836-CC-BY-4.0.txt.\par\smallskip

\noindent\textit{Editorial scope.} Counted unit: the Theorem whose source label is \texttt{t1} in the article (source0.tex lines 358--370, source label \texttt{t1}), delivered together with the article's own complete original proof of it (source0.tex lines 372--466, one \texttt{proof} environment of 95 lines). No mathematics is written, repaired or re-derived: statement and proof are the article's own text line for line. Required context retained uncounted: l1 (lines 295--299). Every definition, display and label of those lines is retained unchanged. Omitted from this package and not redistributed: 1--294, 300--357, 371--371, 467--761. Nothing inside a retained excerpt is omitted or altered: every retained body is delivered exactly as the article's lines are written, so no omission receipt of any kind accompanies this package. Citations: the retained excerpts carry no citation command at all, so no bibliography entry of the article is transcribed and no reference is redistributed. Reference adaptation: none was needed and none was made; every reference inside the delivered unit resolves inside it, so no unresolved reference or citation is produced. Printed slips kept literal: no printed word, symbol, hypothesis or number of the counted statement or of its proof was altered. Layout and class: the article's own class is \texttt{amsart} with packages including graphicx, amssymb, amsthm, amsfonts, latexsym, epsfig, amsmath, amscd, amsbsy, color, tikz, enumerate, epstopdf; the wrapper is the lane's pinned \texttt{amsart} editorial wrapper at 10pt on A4, which restates the theorem-like environments the retained text needs and adds the article's own macro definitions that the retained text needs (\texttt{\textbackslash C}, \texttt{\textbackslash Cl}, \texttt{\textbackslash N}, \texttt{\textbackslash R}, \texttt{\textbackslash eps}, \texttt{\textbackslash im}, \texttt{\textbackslash phi}, \texttt{\textbackslash re}), with extra wrapper packages (bm, bbm, dsfont, xspace, cancel, mathtools). The delivered numbering is the unit's own. Assets: no retained span contains a figure environment, a graphics-inclusion command, a PDF-page inclusion, a table or a TikZ picture; no image file is copied into this package, the package declares no asset dependency and no image file is redistributed with it. Licence: version arXiv:2411.17836v2 of this article is distributed under the Creative Commons Attribution 4.0 International licence, as recorded in the arXiv licence metadata for this exact version and shown on the abstract page https://arxiv.org/abs/2411.17836v2. A full rights notice is delivered at licenses/arxiv-2411.17836-CC-BY-4.0.txt. Characters: every retained excerpt is pure ASCII, so the delivered engine, XeLaTeX with its default Latin Modern text font, reports no missing character.
	\begin{lem}\label{l1}
		Let $v\in\C=\R^2$. Then, as $a\to 0$, the broken line joining
		consecutive points $A^nv$ converges to the logarithmic spiral
		$\phi\mapsto ve^{-t\phi}e^{-i\phi}$ locally uniformly in $t$.
	\end{lem}

	\begin{theorem}\label{t1}
		Let $t_0$ be the root of the equation
		\begin{equation}\label{e1}
			\sqrt{2}e^{-\frac74\pi t} = 1 + t
		\end{equation}
		($t_0 \approx 0.0535502597736068$). There exists a lower 
		semi-continuous function $a(t)$ such that if $t_0 < t < 1$, then for 
		$a < a(t)$ (and $b = 1 - ta$) all even images of the point $Z$ are 
		in the fourth quadrant, and all odd images of $Z$ are in the second 
		quadrant. As a consequence, the unstable manifold $W^u_X$ of the 
		fixed point $X$ is attracted to the periodic orbit $\{ P, P' \}$ of 
		period $2$ and $\Cl W^u_X = W^u_X \cup \{ P, P' \}$.
	\end{theorem}

	\begin{proof}
		Fix $t < 1$. Let us consider the spiral described in Lemma~\ref{l1}, 
		with center at $P$ and starting at $Z$. Note that the points $P$ 
		and $Z$ (as well as $X$, $P'$, $Z^k$, $k \in \N$) go to infinity as 
		$a$ approaches zero. For that reason we will scale these points by 
		the factor $a$. In this case the scaled points converge when $a$ goes 
		to zero, as we will see below. Moreover, the below statements about 
		spirals hold if and only if the analogous statements hold for the 
		scaled points. 
		
		The points $P$ and $Z$, scaled by the factor $a$, have the following 
		forms:
		$$aP = \left(\frac{1+t}{1 + t^2}, \ \frac{(1 - ta)(t-1)}{1 + t^2}\right), 
		\ aZ = \left( \frac{2 + a + \sqrt{a^2 + 4 - 4ta}}{2(1 + t)}, \ 0 \right).
		$$ 
		The slope of the segment between $Z$ and $Z^1$, that is the part of the 
		unstable manifold of $X$ which contains $X$, is
		$(-a - \sqrt{a^2 + 4b})/2 = (-a - \sqrt{a^2 + 4 - 4ta})/2$, so it 
		converges to $-1$ when $a$ goes to $0$. Also, as $a$ goes to zero, the 
		linear maps defining our Lozi map on two half-planes go to the same limit. 
		Therefore, the angle between the segments $[Z^1, Z]$ and $[Z, Z^2]$, as 
		well as between the segments $[aZ^1, aZ]$ and $[aZ, aZ^2]$, goes to zero. 
		Hence, in the limit the tangent to our spiral at $aZ$ has slope $-1$. 	
		
		First, we will show that if $t = t_0$ then the spiral is tangent to the 
		$x$-axis. Thus, if $t > t_0$, it does not intersect the $x$-axis. 
		Similarly, if $t = t_1$, where $t_1$ is the root of the equation
		\begin{equation}\label{e2}
			\sqrt{2} e^{-\frac{5\pi}{4}t} = \frac{(1 + t)^2}{1 - t},
		\end{equation}
		then the spiral is tangent to the $y$-axis, so if $t  > t_1$, it does not
		intersect the $y$-axis. However, $t_1 < t_0$, so if $t > t_0$ then our
		spiral does not intersect any of the coordinate axes.
		
		Now we will make calculations in $\C$. Set
		\[
		p = \lim_{a \to 0}aP = \frac{t + 1}{t^2 + 1} + i\frac{t - 1}{t^2 + 1},
		\ \ \ \ \ z = \lim_{a \to 0}aZ = \frac2{t + 1}.
		\]
		To get the tangent line to the spiral to change from slope $-1$ to slope 
		$0$ (above $p$), we have to go along the spiral by the angle $\frac74\pi$. 
		This means that $t_0$ is the root of the equation
		\begin{equation}\label{e3}
			\im\left(p + (z - p)e^{-t \frac74 \pi}e^{-i \frac74 \pi}\right) = 0.
		\end{equation}
		Similarly, $t_1$ is the root of the equation
		\begin{equation}\label{e4}
			\re\left(p + (z - p)e^{-t \frac54 \pi}e^{-i \frac54 \pi}\right) = 0.
		\end{equation}
		
		Equation~\eqref{e3} can be rewritten as
		\[
		\im\left((z - p)e^{-i \frac74 \pi}\right) = -\im(p)e^{t \frac74 \pi}.
		\]
		We have
		\[
		z - p = \frac{(t - 1)^2}{(t + 1)(t^2 + 1)} - i\frac{t - 1}{t^2 + 1},
		\]
		so
		\[
		\frac{z - p}{\im(p)} = \frac{t - 1}{t + 1} - i.
		\]
		Moreover, $e^{-i \frac74 \pi} = \frac{\sqrt{2}}2(1 + i)$.
		Therefore,~\eqref{e3} becomes
		\[
		\frac{\sqrt{2}}2\im\left(\left(\frac{t-1}{t+1}-i\right)(1+i)\right)
		= -e^{t \frac74 \pi},
		\]
		which is equivalent to~\eqref{e1}.
		
		Similar calculations show that~\eqref{e4} is equivalent to~\eqref{e2}.
		The solution is approximately $t_1\approx 0.05019615992097643$, so
		$t_1<t_0$.
		
		Fix $\tau, \ t_0 < \tau < 1$. As we have already shown, our spiral (with 
		center at $P$ and starting at $Z$) does not intersect the coordinate 
		axes (except at the initial point $Z$). By Lemma~\ref{l1}, there exist 
		$\alpha > 0$ and $\eps > 0$ such that for $a < \alpha$ and 
		$\tau - \eps < t < \tau + \eps$, the broken line joining the consecutive 
		even images of the point $Z$ is so close to the spiral that all even 
		images of the point $Z$ are in the fourth quadrant. Let us denote by 
		$S(\tau)$ the set of all such $\alpha$ for which there exists $\eps > 0$ 
		that satisfies the above property. We define $a(\tau) = \sup S(\tau)$. 
		Recall that the function $f$ is called lower semi-continuous at $x_0$ if 
		for every $y < f(x_0)$ there exists a neighborhood $U$ of $x_0$ such that 
		$f(x) > y$ for all $x \in U$. Clearly, the map $\tau \mapsto a(\tau)$ is 
		lower semi-continuous and if $t_0 < \tau < 1$, then for $a < a(\tau)$ (and 
		$b = 1 - \tau a$) all even images of the point $Z$ are in the fourth 
		quadrant. Consequently, all odd images of $Z$ are in the second quadrant, 
		the unstable manifold $W^u_X$ of the fixed point $X$ (which then consists 
		of a segment in the first quadrant, joins consecutive even images of $Z$ 
		in the fourth quadrant and joins consecutive odd images of $Z$ in the 
		second quadrant) is attracted to the periodic orbit $\{P,P'\}$ of period 
		$2$ and $\Cl W^u_X = W^u_X \cup \{ P, P' \}$.
	\end{proof}

\end{document}
